\chapter{Avaluació i Reducció de Grafs}

La fase d'avaluació constitueix el nucli operatiu de l'intèrpret. Una vegada que el sistema ha verificat la correcció sintàctica i ha inferit satisfactòriament el tipus de l'expressió, l'arbre de sintaxi abstracta (AST) es transforma en una estructura de dades dinàmica capaç de suportar l'avaluació canònica mandrosa (\textit{lazy evaluation}). En aquest capítol es detalla com es realitza la transició de l'AST cap a un graf dirigit, els mecanismes d'adreçament en memòria basats en mètodes d'estat pur, i l'estratègia de reducció per camí esquerre (\textit{spine unwinding}) implementada en els mòduls \texttt{HSDef.hs} i \texttt{HSEval.hs}.

\section{De l'AST al Graf}

A diferència dels intèrprets basats en arbres purs, on cada subexpressió s'avalua de forma independent duplicant el treball en cas de variables compartides, aquest intèrpret empra una tècnica de \textbf{reducció de grafs} (\textit{graph reduction}). La diferència fonamental rau en el fet que un graf dirigit permet que múltiples nodes apuntin a una mateixa adreça de memòria. Això possibilita compartir subexpressions en lloc de copiar-les, establint la base física indispensable per a l'avaluació mandrosa.

La transformació de l'AST cap al graf es troba governada per la funció \texttt{ast2graph} definida al mòdul \texttt{HSEval.hs}. Aquesta funció fa un recorregut descendent de l'AST i va construint seqüencialment els nodes corresponents dins del mapa de memòria global (\textit{Heap}):

\begin{lstlisting}[style=haskellstyle, caption=Funció de transformació de l'AST al graf elemental]
ast2graph :: Expr -> HeapState -> [(String,Addr)] -> (Addr, HeapState)
\end{lstlisting}

La signatura rep l'expressió a processar (\texttt{Expr}), l'estat actual del monticle (\texttt{HeapState}) i un entorn d'adreces associatives (\texttt{[(String,Addr)]}) que mapeja els identificadors textuals de les variables cap a les seves posicions físiques existents en el graf. La funció retorna un parell format per l'adreça de l'arrel de la nova estructura generada (\texttt{Addr}) i l'estat actualitzat del monticle.

El funcionament intern es ramifica segons els constructors definits a \texttt{HSDef.hs}:
\begin{itemize}
    \item \textbf{Valors literals (\texttt{Val v}):} Es genera un node atòmic \texttt{NVal v} i s'insereix en l'adreça lliure actual \texttt{a}, incrementant el comptador global per a futures allocacions.
    \item \textbf{Variables (\texttt{Var s}):} L'intèrpret efectua una cerca a la llista d'associacions \texttt{strAddr}. Si la variable ja existeix (és a dir, ha estat lligada per una lambda superior), es retorna la seva adreça original sense crear cap node nou, la qual cosa consolida formalment el \textbf{compartiment de referències}. Si no es troba, s'assumeix que és una variable global o lliure i s'introdueix un node \texttt{NVar s}.
    \item \textbf{Aplicacions (\texttt{App e1 e2}):} Es resol de manera recursiva la construcció de la branca esquerra (\texttt{e1}) i de la branca dreta (\texttt{e2}). Finalment, s'allotja un node connector \texttt{NApp a1 a2} que enllaça ambdues adreces ordenadament.
    \item \textbf{Abstraccions (\texttt{Lam str e}):} Es modela la unió sintàctica inserint un node \texttt{NLam} que apunta tant a la variable lligada (\texttt{NVar str} a la posició \texttt{a+1}) com al cos transformat (\texttt{a1}), forçant que durant la conversió del cos qualsevol aparició del token quedi associada inequívocament a aquesta posició.
    \item \textbf{Condicionals (\texttt{If e1 e2 e3}):} Es construeix un node de control combinat \texttt{NIf a1 a2 a3} que manté penjats els subgrafs de la condició, de la branca \textit{then} i de la branca \textit{else}.
\end{itemize}

\subsection{Estructura del Heap i Adreçament}

Per modelar el comportament del monticle de memòria conservant la puresa referencial de Haskell, s'ha definit un sistema d'adreçament lògic indexat a \texttt{HSDef.hs}:

\begin{lstlisting}[style=haskellstyle, caption=Definicions de l'estructura del monticle a HSDef.hs]
type Addr = Int
type Heap = IM.IntMap Node
type HeapState = (Addr, Heap)
\end{lstlisting}

L'ús de \texttt{Data.IntMap} (un mapa on les tecles són exclusivament de tipus \texttt{Int}) com a base per representar el graf aporta avantatges arquitectònics molt significatius en comparació amb estructures d'arbres balançats genèrics o llistes de cerca lineals:
\begin{enumerate}
    \item \textbf{Rendiment quasi constant ($O(\min(n, W))$):} Internament, \texttt{IntMap} s'implementa com un arbre de prefixos de Patricia de tipus Big-Endian. Això significa que les operacions de cerca (\texttt{lookup}), inserció (\texttt{insert}) i eliminació (\texttt{delete}) no depenen de forma directa del nombre total de nodes $n$, sinó del nombre de bits $W$ de la clau (que en arquitectures de 64 bits és un màxim fix de 64). A la pràctica, per a la mida habitual dels grafs executats en l'intèrpret, aquest comportament és asimptòticament indistingible d'un rendiment d'ordre constant \textbf{$O(1)$}.
    \item \textbf{Adreçament explícit sense punters reals:} Atès que Haskell és un llenguatge d'alt nivell amb recol·lector de brossa propi, no és possible manipular adreces físiques de memòria RAM de forma directa. L'ús de tipus de dades \texttt{Addr = Int} simula de manera perfecta una memòria d'accés aleatori controlada per programari.
    \item \textbf{Filat d'Estat Pur (\textit{State-Threading}):} L'estat del graf es propaga de manera transparent mitjançant el parell \texttt{HeapState}. El primer element representa la següent adreça lògica lliure disponible (\textit{next-free pointer}), actuant com un comptador incremental que assegura la unicitat de cada node allotjat, eliminant col·lisions d'adreçament sense requerir mutabilitat lateral.
\end{enumerate}

\section{Lambda Lifting i Supercombinadors}

L'avaluació eficient de llenguatges funcionals mitjançant reducció de grafs requereix l'eliminació total de les abstraccions lambda niuades que contenen dependències locals o variables ocultes del context circumdant. El formalisme utilitzat per resoldre aquest inconvenient és el procés de \textbf{Lambda Lifting}, el qual transforma un programa amb definicions de funcions locals en un conjunt homogeni de funcions globals pures conegudes com a \textbf{Supercombinadors}.

A la base de codi de l'intèrpret (\texttt{HSDef.hs}), l'entorn global diferència aquestes categories estructurals mitjançant el tipus algebraic \texttt{EnvObj}:

\begin{lstlisting}[style=haskellstyle, caption=Representació d'objectes de l'entorn global]
data EnvObj = PrimDef Primitive Int    -- Primitiva nuclear i la seva aritat
            | FuncDef Expr            -- Funció pròpia (Supercombinador) i el seu AST
\end{lstlisting}

Un \texttt{FuncDef} emmagatzema l'expressió purificada d'un supercombinador. Quan el motor d'avaluació ensopega amb el nom d'aquesta funció global (un node \texttt{NVar s} registrat a \texttt{DefEnv}), utilitza la plantilla de l'AST desada per instanciar directament un subgraf fresc dins del monticle actiu usant \texttt{ast2graph}, substituint els paràmetres formals per les adreces reals exposades en la pila d'execució.

\subsection{Extracció de variables lliures}

Una variable es diu lliure respecte a una abstracció lambda si apareix dins del cos de la funció però no està lligada pel paràmetre d'aquesta lambda ni per cap entorn global. Per exemple, a l'expressió $\lambda x . (\lambda y . + \ x \ y)$, la variable $x$ és lliure dins de l'abstracció interna que defineix el paràmetre $y$. 

L'extracció de variables lliures és el primer pas del \textit{Lambda Lifting}. Consisteix a analitzar recursivament l'arbre sintàctic per calcular el conjunt de tokens lliures d'una subexpressió. Un cop identificades, aquestes variables s'afegeixen com a paràmetres explícits i formals de la funció interna. En l'exemple anterior, l'abstracció interna es reescriu afegint $x$ com un argument addicional, transformant-se en una estructura equivalent a $(\lambda x . (\text{f\_interna} \ x))$.

\subsection{Generació de funcions globals}

Un cop s'han extret totes les variables lliures d'una funció interna i s'han convertit en paràmetres, la lambda esdevé un \textbf{combinador pur} (manca completament de vincles externs). En aquest punt, es produeix la generació de la funció global: l'abstracció s'extreu completament de la seva posició local dins de l'AST i es promociona com una definició independent al catàleg de primer nivell del sistema (\texttt{preludeDefEnv} o equivalents). 

Aquest procés garanteix que el motor d'avaluació del graf només hagi de gestionar definicions globals planes de tipus \texttt{FuncDef}. Quan s'invoca el supercombinador des del codi, l'intèrpret simplement aplica reduccions posicionals directes sobre la pila de termes, eliminant la necessitat de mantenir tancaments lògics dinàmics clàssics (\textit{closures}) en temps d'execució i optimitzant dràsticament l'ús de la memòria del \texttt{Heap}.

\section{Estratègia de Reducció (Lazy Evaluation)}

L'estratègia d'avaluació de l'intèrpret segueix un model d'\textbf{avaluació mandrosa} o \textit{lazy evaluation}, formalitzada mitjançant la reducció d'ordre normal. Sota aquest model, les expressions no s'avaluen en el moment de ser passades com a arguments d'una funció, sinó únicament i exclusiva quan el seu valor numèric o lògic és estrictament demandat per una operació primitiva primària.

L'estat d'estabilitat d'un node es determina mitjançant el concepte de \textbf{Forma Normal de Cap Feble} o \textbf{WHNF} (\textit{Weak Head Normal Form}). La funció \texttt{isWHNF} codifica aquesta condició de parada:

\begin{lstlisting}[style=haskellstyle, caption=Funció de control de la Forma Normal de Cap Feble]
isWHNF :: (Addr, HeapState) -> Bool
isWHNF (a, (aa,hp)) =
  let Just n = IM.lookup a hp
  in case n of
    NVal _    -> True
    NVar _    -> True
    NLam _ _  -> True
    NApp _ _  -> (
      let (b1, nApps, nPars) = getNumAppOfPredFunc (a, (aa,hp)) 0
      in case b1 of
          True  -> if nApps >= nPars then False else True
          False -> False
      )
    NIf _ _ _ -> False
\end{lstlisting}

Un node es considera estabilitzet en WHNF si és un literal numèric (\texttt{NVal}), una variable lliure sense codi associat (\texttt{NVar}), o una abstracció lambda (\texttt{NLam}), ja que el motor no avalua mai el cos intern d'una lambda fins que aquesta no rebi el seu corresponent argument. En el cas dels nodes d'aplicació (\texttt{NApp}), la funció \texttt{getNumAppOfPredFunc} interroga la branca esquerra per verificar si ens trobem davant d'una primitiva predefinida. Si el nombre d'aplicacions detectades a la pila (\texttt{nApps}) és inferior a l'aritat requerida per la primitiva (\texttt{nPars}), l'expressió representa una \textbf{aplicació parcial} i s'identifica legítimament com a WHNF. En cas contrari, o si es tracta d'un condicional \texttt{NIf}, el bloc requereix més reduccions i retorna \texttt{False}.

\subsection{Exploració del Camí Esquerre (Spine Unwinding)}

El corrent de control de l'avaluació incremental s'executa a la funció \texttt{pas}, la qual rep l'adreça arrel d'estudi, l'estat del monticle, l'entorn de definicions i una llista d'adreces que actua com a \textbf{pila de la línia dorsal} o \textbf{Spine Stack} (\texttt{[Addr]}).

Quan el sistema s'enfronta a una expressió d'aplicació (\texttt{NApp a1 a2}), s'activa el mecanisme conegut com a \textbf{Spine Unwinding} (exploració del camí esquerre). Donat que les aplicacions de funcions són associatives cap a l'esquerra, la funció real es troba amagada a l'extrem inferior esquerre de la cadena de nodes. El sistema realitza un descens recursiu acumulant l'adreça de cada node d'aplicació directament a la pila estructural:

\begin{lstlisting}[style=haskellstyle, caption=Cas de l'Spine Unwinding a la funció pas]
NApp a1 a2 -> pas (a1, (aa,hp)) defEnv (a:stk)
\end{lstlisting}

Aquest descens continua de forma ininterrompuda fins a trobar un node que no sigui una aplicació (normalment una lambda \texttt{NLam} o una variable de primitiva \texttt{NVar}). En aquest precís instant, la pila (\texttt{stk}) conté, ordenades de dalt a baix, les adreces dels nodes \texttt{NApp} que capturen els arguments de la funció.

Un cop assolit el fons esquerre, es bifurca la reducció segons el node descobert:
\begin{itemize}
    \item \textbf{Beta-reducció (\texttt{NLam a1 a2}):} Es recupera el node d'aplicació superior de la pila (\texttt{app1}), el qual conté l'argument real a la seva branca dreta (\texttt{a3}). S'extreu el contingut d'aquest argument i s'enllaça directament substituint la variable formal \texttt{a1}. Posteriorment, es realitza un pas crucial d'actualització: el node d'aplicació \texttt{app1} es \textbf{reescriu completament en memòria} col·locant-hi el contingut del cos de la lambda (\texttt{brancaDretaLam}). Això garanteix que qualsevol altra part del graf que compartís aquest node d'aplicació visualitzi de forma instantània el resultat de la reducció, evitant duplicació de càlculs posteriors i consolidant el comportament mandrós.
    \item \textbf{Avaluació de Primitives (\texttt{NVar s} apuntant a un \texttt{PrimDef}):} S'extreuen de la pila els nodes d'aplicació necessaris segons l'aritat de l'operador. Atès que les primitives d'aquest intèrpret (com \texttt{Plus} o \texttt{Mult}) són \textbf{estrictes}, es requereix forçar l'avaluació prèvia dels seus arguments abans de realitzar el càlcul matemàtic. La funció comprova si els fills estan en WHNF; si no ho estan, crida recursivament a \texttt{pas} focalitzant-se en la posició de l'argument desprotegit. Un cop ambdós operands han col·lapsat cap a nodes \texttt{NVal}, s'invoca la funció de càlcul \texttt{compute}, es converteix el resultat de nou a subgraf a través d'\texttt{ast2graph}, i s'actualitza el node d'aplicació arrel de l'operació per contenir el valor final.
\end{itemize}

Aquest disseny d'actualitzacions \textit{in situ} sobre el \texttt{IntMap} d'estat pur garanteix l'harmonia absoluta entre la immutabilitat lògica del llenguatge de programació Haskell i l'eficiència pràctica d'una màquina de reducció de grafs d'alt rendiment. El funcionament es pot visualitzar clarament mitjançant la reducció de l'expressió guia \texttt{(\textbackslash x y -> x * y) 2 3}: en entrar al bloc \texttt{debugEvalLoop}, l'arbre construït per \texttt{ast2graph} s'explora seguint l'espinada esquerra fins a localitzar la lambda, moment en què s'extreuen seqüencialment les adreces dels literals \texttt{2} i \texttt{3} de la pila \texttt{stk}, es reescriuen els nodes intermedis i finalment l'operació \texttt{Mult} de la primitiva unifica el resultat en un node unificat \texttt{NVal 6}.
